\section{Introduction}
\label{sec:introduction}
An idea to minimize a difficult target objective $L(\theta)$ is to introduce an approximation family $(L_s)_{s\in[0,1]}$, with $L_1=L$, such that the initial problems are easier to solve. The family progressively evolves to the target problem. In practice, one considers $0=s_0<\cdots<s_K=1$ and uses the solution obtained at one stage as a warm start for the next. In the best case, the family admits a path of minimizers that a local optimizer can follow. Graduated optimization gives guarantees of this kind under specific geometric assumptions \citep{hazan2016graduated}. Whether a useful family can be constructed for neural-network training is less clear.

Our starting point was to explore several such families. Directly deforming the objective, adding noise in parameter space, or modifying the activations all provide possible approaches. However, choosing an easier initial problem is already difficult, and explicitly smoothing a high-dimensional loss can be expensive. We therefore retained a more concrete starting point: the Gaussian filtering of CNN feature maps proposed by \citet{sinha2020curriculum}.

For image-based learning, the intuition is that the model can first focus on the large-scale structures of the images before having to deal with finer details and textures. We then viewed the Gaussian kernel more generally as a way of controlling information, and considered other approaches of doing the same thing. This led us from filtering to temporary resolution reduction, including reduction inside the network. The question became: how much of the observed benefit comes from Gaussian smoothing, and how much can be obtained by simply processing smaller feature maps early in training?

Our exploratory experiments suggest that a simple internal max-pooling intervention already accounts for a substantial part of the gain. In the final batch, where all procedures share initial weights and data order seed by seed, reduction after the second residual block improves test accuracy by $\uGainRes$ percentage points over the baseline on average. Adding Gaussian filtering at convolution outputs increases this gain to $\uGainComb$ points. These results hold for the tested short CIFAR-10 training protocol, on a test set that was also used to select the procedures; they do not establish an improvement over a fully tuned training recipe. The detailed exploration is reported in Appendix~\ref{app:exploration}.

We now restrict the study to three fixed methods: resolution reduction, Gaussian filtering, and their combination. Our aim is to establish how well these methods transfer and to understand what they change during optimization. In particular, a lower training loss is not necessary for a better test loss in our experiments: at the end of training, the three methods have a higher cross-entropy than the baseline on a fixed subset of training images, and a lower test cross-entropy. This motivates both a statistical interpretation and an examination of the loss along the training trajectories, rather than another large search over operators.
